\chapter{Dataflow and Apache Beam Fundamentals}
\index{Dataflow}\index{Apache Beam}\index{PCollection}\index{PTransform}\index{ParDo}\index{Flex Templates}

\begin{objectives}
\begin{itemize}
    \item Apply the Beam programming model: pipelines, PCollections, PTransforms, ParDo, and runners.
    \item Explain Dataflow's managed architecture: worker autoscaling, Shuffle and Streaming Engine.
    \item Unit-test pipelines with TestPipeline, PAssert, and isolated DoFn tests before deploying.
    \item Package pipelines as Classic or Flex Templates for repeatable, parameterized launches.
    \item Build the WordCount pipeline and run it on the Dataflow runner with autoscaling.
    \item Design a real-time pipeline that anonymizes healthcare HL7 messages before warehousing.
\end{itemize}
\end{objectives}

\begin{crosscloud}
Managed runners for portable pipelines exist on every cloud, but Beam's model---one API, many runners---is the most explicit separation of pipeline logic from execution engine.

\vspace{2mm}
{\footnotesize
\begin{tabular}{@{}p{2.8cm}p{3.0cm}p{3.0cm}p{3.0cm}@{}} \\
\toprule
\textbf{Capability} & \textbf{GCP Dataflow} & \textbf{AWS} & \textbf{Azure} \\
\midrule
Managed runner & Dataflow (Beam native) & Kinesis Data Analytics (Flink), EMR Flink/Spark & HDInsight/Synapse Spark, Fabric \\
Programming model & Beam (Java/Python/Go) & Flink API, Spark API, Kinesis SQL & Spark API, KQL \\
Portability & Beam runners (Flink, Spark, Samza) & Moderate & Moderate \\
Autoscaling & Native, per-stage & Application scaling & Pool-level scaling \\
Batch + streaming & One model, both modes & Separate services traditionally & Separate engines traditionally \\
\bottomrule
\end{tabular}}

\vspace{1mm}
The Beam model's contribution is conceptual as much as practical: its unified treatment of bounded and unbounded data, event time, and windowing \cite{akidau2015dataflow} is now the vocabulary in which streaming is discussed on every platform, including \textbf{Snowflake}'s streaming features and \textbf{MongoDB}'s change-stream processors.
\end{crosscloud}

\begin{keyterms}
\textbf{PCollection}: an immutable, distributed dataset---bounded or unbounded---flowing through a pipeline.
\textbf{PTransform}: a data-processing operation consuming and producing PCollections.
\textbf{ParDo}: the general parallel ``do'' transform executing a user DoFn per element.
\textbf{Runner}: the engine executing a Beam pipeline---Direct (local), Dataflow, Flink, Spark.
\textbf{Streaming Engine}: Dataflow's mode that moves streaming shuffle and state off worker VMs into the service backend.
\textbf{Flex Template}: a containerized pipeline image plus a template spec, launched with parameters without rebuilding code.
\end{keyterms}

\section{Week 14 Roadmap}
Week~14 teaches the model and the managed runner together. Monday is the Beam model: PCollections, PTransforms, ParDo, and the runner abstraction. Tuesday is Dataflow's service architecture: autoscaling, Shuffle, Streaming Engine. Wednesday covers templates. Thursday runs WordCount on the Dataflow runner. Friday anonymizes healthcare HL7 streams before they reach BigQuery.

\begin{definitionbox}
\textbf{Apache Beam} is a unified programming model for batch and streaming data processing. A Beam pipeline is a graph of PTransforms over PCollections, compiled to a \textbf{runner}; \textbf{Dataflow} is Google Cloud's fully managed Beam runner with autoscaling, exactly-once shuffle, and streaming state management \cite{apachebeamdocs,googleclouddataflowdocs}.
\end{definitionbox}

\section{Monday: The Beam Model---PCollections, PTransforms, ParDo, Runners}
\index{Beam model}\index{DoFn}\index{runners}

\subsection{Pipelines as Graphs}
A Beam program builds a directed graph: sources produce PCollections; transforms map, filter, key, window, group, and join them; sinks write results. The same graph runs bounded (a batch of files) or unbounded (a Pub/Sub subscription)---the model unifies both, which is the thesis of the Dataflow paper \cite{akidau2015dataflow}: ask \emph{what} results, \emph{where} in event time, \emph{when} in processing time to emit, and \emph{how} refinements accumulate, and the model handles the rest.

\subsection{ParDo and DoFn}
\texttt{ParDo} is the workhorse: your \texttt{DoFn.process} runs per element, emitting zero or more outputs, with access to side inputs, state, and timers (Chapter~15). Higher-level transforms---\texttt{Map}, \texttt{Filter}, \texttt{GroupByKey}, \texttt{CombinePerKey}---are convenience wrappers over the same machinery. Fusions and combiner lifting are why idiomatic Beam (combine before group, filter early) runs dramatically cheaper than naive graphs.

\begin{lstlisting}[language=Python,caption={The canonical WordCount expressed in Beam.}]
import re
import apache_beam as beam

with beam.Pipeline() as p:
    (p
     | "Read" >> beam.io.ReadFromText("gs://dataflow-samples/shakespeare/kinglear.txt")
     | "Words" >> beam.FlatMap(lambda line: re.findall(r"[A-Za-z']+", line))
     | "Pair" >> beam.Map(lambda w: (w.lower(), 1))
     | "Count" >> beam.CombinePerKey(sum)
     | "Format" >> beam.Map(lambda kv: f"{kv[0]}: {kv[1]}")
     | "Write" >> beam.io.WriteToText("gs://my-bucket/out/wordcount"))
\end{lstlisting}

\subsection{Runners and Portability}
Develop on the Direct runner locally, test with TestPipeline, deploy to Dataflow. Because the model is portable, the same pipeline can run on Flink or Spark---a hedge some organizations value---but portability is at the API level: runner-specific features (Streaming Engine, exactly-once modes) are where managed value concentrates.

\begin{notebox}
Beam pipelines are lazily constructed: transforms build the graph, and \texttt{pipeline.run()} (or context exit) submits it. Code before submission runs on the launching machine; DoFn bodies run on workers. Confusing the two is the most common beginner bug---a database connection opened at graph-construction time does not exist on workers.
\end{notebox}

\section{Tuesday: Dataflow Architecture---Autoscaling, Shuffle, Streaming Engine}
\index{Dataflow!autoscaling}\index{Streaming Engine}\index{shuffle}

\subsection{The Managed Service}
You submit a graph; Dataflow provisions worker VMs, distributes fused stages, and scales workers between configured bounds as backlog and utilization demand. The service monitors per-stage progress, drains stragglers with dynamic work rebalancing (stealing unprocessed ranges from slow workers), and reports per-stage metrics---system lag, watermark, throughput---that become your alerting surface in Chapter~15.

\subsection{Shuffle and Streaming Engine}
Batch jobs use the service-based \textbf{Shuffle}: intermediate data moves through a managed shuffle tier instead of worker disks, so workers stay stateless and autoscale freely. \textbf{Streaming Engine} does the same for streaming pipelines---plus moving per-key state into the backend---reducing worker count, improving autoscaling latency, and enabling features like exactly-once shuffle. Unless you have a specific reason (custom network placement, specialized disks), both should be on.

\begin{table}[H]
\centering
\small
\begin{tabular}{@{}p{4.4cm}p{4.4cm}p{4.4cm}@{}} \\
\toprule
\textbf{Signal} & \textbf{Meaning} & \textbf{Response} \\
\midrule
Workers at max, backlog growing & Under-provisioned or hot stage & Raise \texttt{max\_num\_workers}; inspect stage \\
Workers idle, low throughput & Input-bound or over-provisioned & Check source parallelism; lower bounds \\
High system lag, stable watermark & Processing capacity short & Scale up; enable Streaming Engine \\
Watermark stuck, lag stable & Late/stranded data upstream & Inspect upstream; check event timestamps \\
\bottomrule
\end{tabular}
\caption{Reading Dataflow job metrics during tuning.}
\label{tab:ch14-dataflow-metrics}
\end{table}

\begin{tipbox}
Autoscaling reacts to backlog; it does not predict it. For known daily peaks, either accept the scale-up lag (minutes) or pin a higher \texttt{num\_workers} floor during the peak window via scheduled job updates.
\end{tipbox}

\section{Wednesday: Templates---Classic versus Flex}
\index{templates!classic}\index{templates!flex}

Templates separate pipeline \emph{definition} from pipeline \emph{launch}: operators launch parameterized jobs without a development environment. \textbf{Classic templates} stage the serialized graph at build time---parameters must be \texttt{ValueProvider}s resolved at launch, and graph shape cannot depend on them. \textbf{Flex Templates} package the pipeline as a container image; at launch the container runs, builds the graph with full access to parameters, and submits it---conditional graph construction, custom dependencies, and non-ValueProvider parameters all work. Google ships both as Google-provided templates (GCS to BigQuery, Pub/Sub to BigQuery, and dozens more); prefer them over writing your own when they fit.

\begin{pitfall}
A Classic template whose graph depends on a launch parameter silently bakes the build-time value. If a parameter must change the DAG shape---choose a source, toggle a transform---you need a Flex Template. This constraint has bitten every team that ignored it exactly once.
\end{pitfall}

\section{Thursday Hands-On Project: WordCount on the Dataflow Runner}
\index{hands-on project}\index{WordCount}\index{unit testing!Beam}
This lab runs WordCount locally, unit-tests a transform, then submits to Dataflow with autoscaling---the full develop-test-deploy loop in miniature.

\subsection{Local Run and Unit Test}
\begin{lstlisting}[language=Python,caption={Unit-testing a Beam transform with TestPipeline and PAssert.}]
import apache_beam as beam
from apache_beam.testing.test_pipeline import TestPipeline
from apache_beam.testing.util import assert_that, equal_to


def count_words(lines):
    return (lines
            | "Words" >> beam.FlatMap(str.split)
            | "Pair" >> beam.Map(lambda w: (w, 1))
            | "Count" >> beam.CombinePerKey(sum))


def test_count_words():
    with TestPipeline() as p:
        out = p | beam.Create(["a b a", "b"])
        assert_that(count_words(out), equal_to([("a", 2), ("b", 2)]))
\end{lstlisting}

\subsection{Submit to Dataflow}
\begin{lstlisting}[language=bash,caption={Running WordCount on Dataflow with autoscaling from PowerShell.}]
$env:PROJECT_ID = "my-gcp-dataflow-lab"
$env:REGION = "us-central1"
$env:BUCKET = "my-gcp-dataflow-lab-bucket"

python .\wordcount.py `
  --runner=DataflowRunner `
  --project=$env:PROJECT_ID `
  --region=$env:REGION `
  --temp_location=gs://$env:BUCKET/tmp `
  --staging_location=gs://$env:BUCKET/staging `
  --output=gs://$env:BUCKET/out/wordcount `
  --autoscaling_algorithm=THROUGHPUT_BASED `
  --max_num_workers=4
\end{lstlisting}

Open the Dataflow job graph in the console: identify the fused stages, worker count over time, and per-stage throughput. Then rerun with \texttt{--runner=DirectRunner} on a small input and compare---the graph is identical; only the runner changed.

\subsection*{Deliverables}
\begin{itemize}
    \item The pipeline, the passing unit test, and the submission command transcript.
    \item A screenshot of the Dataflow job graph with stages and autoscaling behavior annotated.
    \item A paragraph comparing the local and managed runs: what the runner changed and what it did not.
\end{itemize}

\subsection*{Acceptance Criteria}
\begin{itemize}
    \item The transform logic is covered by a PAssert unit test that fails when logic changes.
    \item The Dataflow job succeeds with autoscaling enabled and bounded workers.
    \item Outputs are written to a deterministic Cloud Storage path, overwritable on rerun.
\end{itemize}

\subsection*{Stretch Goals}
\begin{itemize}
    \item Package the pipeline as a Flex Template with a Docker image and launch it with \texttt{gcloud dataflow flex-template run}.
    \item Add a Google-provided template run (Text Files on Cloud Storage to BigQuery) and compare with hand-rolled code.
    \item Benchmark the job with and without service-based shuffle on a large input.
\end{itemize}

\section{Friday Use Case and Architecture: Real-Time HL7 Anonymization}
\index{HL7}\index{healthcare anonymization}\index{de-identification!pipeline}

\subsection{Scenario}
A hospital network streams HL7 admission and lab messages from facilities to a central analytics platform. Before any message lands in BigQuery, direct identifiers (patient names, MRNs, addresses) must be pseudonymized; the mapping from real to pseudonymous identity lives in a restricted system; analysts query only the safe projection.

\subsection{Architecture}
Facilities publish to a schema-validated Pub/Sub topic (Chapter~13). A streaming Dataflow pipeline parses HL7 segments, calls the tokenization service (or Sensitive Data Protection's de-identification, Chapter~23) through a bounded-batch \texttt{ParDo} with backoff, and writes two outputs: the pseudonymized stream to BigQuery for analytics, and the raw stream to a VPC-SC-protected, CMEK-encrypted Cloud Storage bucket reachable only by the compliance service account. Side inputs carry the day's code mappings; dead-lettered parse failures go to a triage topic with alerts.

\begin{figure}[H]
    \centering
    \resizebox{0.95\textwidth}{!}{%
    \begin{tikzpicture}[
        src/.style={rectangle, rounded corners, draw=codegray, thick, fill=white, text width=2.7cm, minimum height=1.0cm, align=center, font=\small\bfseries},
        svc/.style={rectangle, rounded corners, draw=primary, thick, fill=primary!6, text width=3.0cm, minimum height=1.0cm, align=center, font=\small\bfseries},
        dst/.style={rectangle, rounded corners, draw=codegreen!60!black, thick, fill=green!7, text width=2.9cm, minimum height=1.0cm, align=center, font=\small\bfseries},
        warn/.style={rectangle, rounded corners, draw=red!60!black, thick, fill=red!6, text width=2.9cm, minimum height=0.9cm, align=center, font=\footnotesize\bfseries},
        arrow/.style={-{Stealth[length=2.5mm]}, draw=primary, thick},
        dasharrow/.style={-{Stealth[length=2.5mm]}, draw=codegray, thick, dashed}
    ]
        \node[src] (hl7) at (0,0) {Facilities\\HL7 feeds};
        \node[svc] (ps) at (3.6,0) {Pub/Sub\\schema topic};
        \node[svc] (df) at (7.4,0) {Dataflow\\parse + tokenize\\(batch API calls)};
        \node[dst] (bq) at (11.6,1.2) {BigQuery\\pseudonymous\\analytics};
        \node[warn] (raw) at (11.6,-1.2) {GCS raw vault\\VPC-SC + CMEK};

        \draw[arrow] (hl7) -- (ps);
        \draw[arrow] (ps) -- (df);
        \draw[arrow] (df) -- (bq);
        \draw[dasharrow] (df) -- (raw);
    \end{tikzpicture}%
    }
    \caption{HL7 anonymization: parse and tokenize in flight; analysts see only the pseudonymous projection.}
    \label{fig:ch14-hl7-anonymization}
\end{figure}

\begin{tipbox}
Batch external API calls inside the DoFn (accumulate a bundle, call once) and bound concurrency with backoff---per-message calls to a tokenization endpoint are the classic way to turn a cheap pipeline into an expensive, rate-limited one. Chapter~20 develops this pattern fully.
\end{tipbox}

\begin{pitfall}
``Anonymized'' pipelines leak through timestamps, free-text notes, and rare-diagnosis combinations. Pseudonymization of direct identifiers is step one; quasi-identifier risk analysis (Chapter~23) decides whether the analytics projection is actually safe to expose.
\end{pitfall}

\section{Interview Q\&A}
\subsection{Concepts and Trade-offs}
\begin{enumerate}
    \item \textbf{Q: What is the difference between a PTransform and a ParDo?}\\
    A: A PTransform is any graph node transforming PCollections; ParDo is the specific general-purpose transform applying your DoFn element-wise.

    \item \textbf{Q: Why unit-test Beam with TestPipeline before deploying to Dataflow?}\\
    A: Logic errors cost minutes locally and money plus data-correction effort in production; PAssert proves transform semantics, and TestStream (Chapter~15) proves windowing behavior, all in CI.

    \item \textbf{Q: What does Streaming Engine move off the worker VMs?}\\
    A: Streaming shuffle and per-key state into the service backend, leaving workers stateless, smaller, and faster to autoscale.

    \item \textbf{Q: When should you package a pipeline as a Flex Template?}\\
    A: When operators must launch parameterized jobs without a dev environment, or when graph shape depends on launch parameters---which Classic templates cannot express.

    \item \textbf{Q: How does Dataflow autoscaling react to a growing backlog?}\\
    A: It adds workers up to \texttt{max\_num\_workers} based on backlog and utilization, then removes them as the backlog drains---reactive, with minutes of ramp time.
\end{enumerate}

\section{Further Reading}
Read the Beam programming guide for the model, transforms, and testing utilities \cite{apachebeamdocs}, and the Dataflow documentation for runner options, Streaming Engine, and templates \cite{googleclouddataflowdocs}. The Dataflow Model paper grounds the unified batch/streaming semantics \cite{akidau2015dataflow}. Chapter~13's Pub/Sub documentation covers the ingestion contracts feeding these pipelines \cite{googlecloudpubsubdocs}.

\section*{Key Takeaways}
\begin{itemize}
    \item Beam unifies batch and streaming: one graph, bounded or unbounded inputs, many runners.
    \item ParDo/DoFn is where user logic lives; prefer combinable transforms and early filters for cheap graphs.
    \item Streaming Engine and service shuffle make workers stateless---enable both by default.
    \item Autoscaling reacts to backlog within worker bounds; plan peaks deliberately.
    \item Test with TestPipeline and PAssert before any managed run; templates industrialize launches.
    \item Regulated pipelines separate raw vaults from pseudonymous analytics at the architecture level, not by convention.
\end{itemize}

\section{Exercises}
\begin{enumerate}
    \item \textbf{(Conceptual)} Explain graph-construction time versus worker execution time, with one bug that arises from confusing them.
    \item \textbf{(Conceptual)} Why does moving state into Streaming Engine improve autoscaling for streaming pipelines?
    \item \textbf{(Hands-on)} Reproduce the Thursday lab; then modify the tokenizer to strip stopwords, update the PAssert, and redeploy.
    \item \textbf{(Hands-on)} Package WordCount as a Flex Template and launch two runs with different parameters.
    \item \textbf{(Analysis)} A streaming job shows rising system lag with a stable watermark. Diagnose using Table~\ref{tab:ch14-dataflow-metrics} and propose the fix.
    \item \textbf{(Design)} Extend the HL7 design with a Bigtable latest-state store per patient pseudonym, including row-key sketch and TTL policy.
\end{enumerate}
